\chapter{交换代数. 代数数论.}

“抽象”交换代数是一项晚近的创造，但若不联系孕育了它的代数数论与代数几何的发展，其发展过程便无从理解。

不大离谱地可以猜测：费马自称掌握的那个著名“证明”——方程 \(x^{p} + y^{p} = z^{p}\)（\(p\) 为奇素数，\(x, y, z\) 为整数 \(\neq 0\)）不可能成立——乃是基于分解式

\[
(x + y) (x + \zeta y) \dots (x + \zeta^ {p - 1} y) = z ^ {p}
\]

在环 \(\mathbf{Z}[\zeta]\)（其中 \(\zeta \neq 1\) 是一个 \(p\) 次单位根）中的分解，并且基于（在假定该环为主理想环的情况下）关于这个环中整除性的一个论证。无论如何，人们发现拉格朗日已经概述过一个类似的论证（{[}191{]}，第 II 卷，第 531 页）；正是依靠这类论证及其各种变体（特别是以降低方程次数为目的的变量代换），欧拉（{[}108 a{]}，第 I 卷，第 488 页）\(^{1}\)和高斯（{[}124 a{]}，第 II 卷，第 387 页）证明了 \(p = 3\) 时的费马定理，高斯（同前引文）与狄利克雷（{[}92{]}，第 I 卷，第 42 页）证明了 \(p = 5\) 的情形，狄利克雷还证明了方程 \(x^{14} + y^{14} = z^{14}\) 的不可解性（{[}92{]}，第 I 卷，第 190 页）。最后，库默尔在其关于数论的最初研究中，曾相信自己已用这种方法得到了一个一般的证明，而无疑正是这个错误（由狄利克雷向他指出）引导他去研究分圆域的算术，并从那里他终于成功地导出其证明的一个正确版本，适用于素数 \(p < 100\) {[}188 d{]}。

另一方面，高斯 1831 年关于双二次剩余的著名论文——其结果出自对“高斯整数”环 Z{[}i{]} 中整除性的细致研究（{[}124 a{]}，第 II 卷，第 109 页）——表明，把整除性的概念推广到代数数，对数论的经典问题会有何等益处；\(^{2}\)因此毫不奇怪，1830 到 1850 年间，这一理论成了众多

德国数学家竞相研究的对象：先是雅可比、狄利克雷和艾森斯坦，稍后一点是库默尔以及他的学生兼朋友克罗内克。我们在这里不打算谈单位数理论，它过于局限于数论，那里进展迅速：艾森斯坦得到了三次域单位数群的结构，克罗内克得到了分圆域的情形，此后不久，狄利克雷于 1846 年（{[}92{]}，第 I 卷，第 640 页）证明了一般的定理，而埃尔米特凭一己之力已几乎达到这个定理（{[}159{]}，第 I 卷，第 159 页）。看来艰难得多的是（整个理论的中心）问题，即分解为素因子的问题。自从拉格朗日给出形如 \(x^{2} + Dy^{2}\)（\(x, y, D\) 为整数）的数具有形如 \(m^{2} + Dn^{2}\) 之外的因子的例子（{[}191{]}，第 II 卷，第 465 页）以来，人们在实质上已经知道，不能指望环 \(\mathbf{Z}[\sqrt{-D}]\) 一般而言是主理想环，于是在欧拉的胆大妄为之后继之以高度的审慎；例如，当狄利克雷证明关系式 \(p^{2} - 5q^{2} = r^{5}\)（\(p, q, r\) 为整数）等价于 \(p + q\sqrt{5} = (x + y\sqrt{5})^{5}\)（\(x, y\) 为整数）时，他只限于在一条注记中指出：“对许多其他素数{[}除 5 以外{]}也有类似的定理”（{[}92{]}，第 I 卷，第 31 页）。有了高斯 1831 年的论文和艾森斯坦关于三次剩余的工作 {[}102 a{]}，人们确实对主理想环 \(\mathbf{Z}[i]\) 与 \(\mathbf{Z}[\rho]\)（\(\rho = (-1 + i\sqrt{3})/2\)，一个三次单位根）的算术有了一些深入的研究，与有理整数的理论完全类似；至少凭借这些例子，二次域中的算术与高斯所发展的二元二次型理论之间的紧密联系显得十分清楚；但对一般的情形，还缺少一部“词典”，使我们能通过简单翻译高斯的理论来处理二次域。\(^{3}\)

事实上，这个谜首先不是在二次域、而恰恰是在分圆域上（而且理由要到很久以后才显得清楚（参见第 98 页））被解开的。从 1837 年起，起初身为分析学家的库默尔转向分圆域的算术，在随后的 25 年里，这方面的工作几乎一直占据着他的全部精力。与他的前辈们一样，他研究环 Z{[}ζ{]}（其中 ζ 是一个 ≠ 1 的 p 次单位根（p 为奇素数））中的整除性；他很快就看出，在那里同样会出现并非主理想环的环，从而阻碍了把算术定律加以推广的一切进展 {[}188 a{]}；直到 1845 年，在 8 年努力之后，光明终于降临，这要归功于他对“理想数”的定义 {[}188 c 与 d{]}。

用现代的语言说，库默尔所做的事情归结为在域 \(\mathbf{Q}[\zeta]\) 上定义赋值：这些赋值与他的“素理想数”成一一对应，一个这样的因子出现在某个数 \(x \in \mathbf{Z}[\zeta]\) 的“分解”式中的“指数”，不是别的，正是 \(x\) 在相应赋值下的值。由于 \(x\) 的共轭元也属于 \(\mathbf{Z}[\zeta]\)，而且它们的乘积 \(N(x)\)（即 \(x\) 的“范数”\(^{4}\)）是有理整数，有待定义的“素理想因子”也必须是有理素数的“因子”，而要给出它的定义，只需说它们是某个素数 \(q \in \mathbf{Z}\) 的“素理想因子”即可。当 \(q = p\) 时，库默尔在实质上已经证明 {[}188 a{]}：主理想 \((1 - \zeta)\) 是素理想，而且它的 \((p - 1)\) 次幂就是主理想 \((p)\)；因此这个情形并未提出任何新问题。当 \(q \neq p\) 时，看来指导了库默尔的想法是：把分圆方程 \(\Phi_p(z) = 0\) 换成同余式 \(\Phi_p(u) \equiv 0 \pmod{q}\)，换句话说，把分圆多项式 \(\Phi_p(X)\) 在域 \(\mathbf{F}_q\) 上分解，并把这个多项式的每个不可约因子与一个“素理想因子”对应起来。一个简单的情形（库默尔在宣布其结果而不给证明的短文 {[}188 b{]} 中明确引用了它）是 \(q \equiv 1 \pmod{p}\) 的情形；若 \(q = mp + 1\)，且 \(\gamma \in \mathbf{F}_q\) 是 1 的一个 \((q - 1)\) 次原根，那么在 \(\mathbf{F}_q[X]\) 中有

\[
\Phi_ {p} (X) = \prod_ {k = 1} ^ {p - 1} (X - \gamma^ {k m})
\]

由于 \(\gamma^{pm} = 1\)。于是，对每个因子 \(X - \gamma^{km}\)，让一个“素理想因子”\(\mathfrak{q}_k\)（\(q\) 的因子）与之相对应，库默尔说：一个元素 \(x \in \mathbf{Z}[\zeta]\)（设其极小多项式为 \(P\)，定义在 \(\mathbf{Q}\) 上）被 \(\mathfrak{q}_k\) 整除，如果在 \(\mathbf{F}_q\) 中有 \(P(\gamma^{km}) = 0\)；总而言之，用现代的语言，他把商环 \(\mathbf{Z}[\zeta]/q\mathbf{Z}[\zeta]\) 写成一些同构于 \(\mathbf{F}_q\) 的域的直积。当 \(q \not\equiv 1 \pmod{p}\) 时，\(\Phi_p(X)\) 在 \(\mathbf{F}_q[X]\) 中的不可约因子不再是一次的，因此必须把 \(X\)（在 \(P(X)\) 中）代之以 \(\Phi_p\) 在 \(\mathbf{F}_q[X]\) 中的因子的“伽罗瓦虚根”。库默尔通过——用今天的话说——过渡到分解域 \(K\)（就 \(q\) 而言）来避开这个困难：若 \(f\) 是使 \(q^f \equiv 1 \pmod{p}\) 成立的最小整数，并令 \(p-1 = ef\)，则 \(K\) 不是别的，正是 \(\mathbf{Q}(\zeta)\) 的这样一个子域：它由某个阶为 \(f\) 的子群的不变量构造而成，该子群包含于一个 \(p-1\) 阶的循环群，即 \(\mathbf{Q}(\zeta)\) 在 \(\mathbf{Q}\) 上的伽罗瓦群；换句话说，它是 \(\mathbf{Q}(\zeta)\) 的唯一的次数为 \(e\)（在 \(\mathbf{Q}\) 上）的子域；它自高斯的《Disquisitiones》以来便为人所熟知，由“周期”

\[
\eta_ {k} = \zeta_ {k} + \zeta_ {k + e} + \zeta_ {k + 2 e} + \dots + \zeta_ {k + (f - 1) e}
\]

\((0 \leq k \leq e - 1, \zeta_{\nu} = \zeta^{g^{\nu}}\)，其中 \(g\) 是同余式 \(z^{p-1} \equiv 1 \pmod{p}\) 的一个原根），它们构成此域的一个正规基。若 \(R(X)\) 是这些“周期”中任意一个 \(\eta\) 的极小多项式（首一、以有理整数为系数），那么库默尔依据高斯的公式证明：在域 \(\mathbf{F}_{q}\) 上，\(R(X)\) 仍分解成互异的一次因子 \(X - u_{j}\)（\(1 \leq j \leq e\)），而这次他是把每个 \(u_{j}\) 与一个“素理想因子”\(q_{j}\) 联系起来。为了定义“被 \(q_{j}\) 整除”，库默尔把每个 \(x \in \mathbb{Z}[\zeta]\) 写成形式 \(x = \sum_{k=0}^{f-1} \zeta^{k} y_{k}\)，其中每个 \(y_{k} \in K\) 本身都以唯一方式写成一个次数 \(\leq e-1\) 的 \(\eta\) 的多项式，系数是有理整数；他说 \(x\) 被 \(q_{j}\) 整除，当且仅当以 \(u_{j}\) 代换 \(\eta\)（在每个 \(y_{k}\) 中都这样做）后，所得到的 \(\mathbf{F}_{q}\) 中的元素全都为零。但还必须定义 \(q_{j}\) 在 \(x\) 中的“指数”。为此，库默尔引入了我们现在会称之为 \(q_{j}\) 的一个一致化元的东西，也就是说，一个元素 \(\rho_{j} \in K\)，使得 \(N(\rho_{j}) \equiv 0 \pmod{q}\)、\(N(\rho_{j}) \not\equiv 0 \pmod{q^{2}}\)，最后还使得 \(\rho_{j}\) 被 \(q_{j}\)（按上面定义的意义）整除，而不被任何别的理想因子 \(\neq q_{j}\)（它们都是 \(q\) 的因子）整除。这样一个 \(\rho_{j}\) 的存在性，已由克罗内克前一年在他的学位论文中于实质上证明（{[}186 a{]}，第 I 卷，第 23 页）；随后令 \(\rho_{j}' = N(\rho_{j}) / \rho_{j}\)，库默尔说，\(q_{j}\) 在 \(x\) 中的指数等于 \(h\)，如果有 \(x\rho_{j}'^{h} \equiv 0 \pmod{q^{h}}\)，但 \(x\rho_{j}'^{h+1} \not\equiv 0 \pmod{q^{h+1}}\)；他当然首先证明：关系 \(x\rho_{j}' \equiv 0 \pmod{q}\) 等价于 \(x\) 被 \(q_{j}\) 整除这一事实（按前面的意义）。这些定义一经建立，对“理想数”而言的通常整除定律向 \(Z[\zeta]\) 的推广便没有带来任何严重的困难；而且从他的第一篇论文 {[}188 a{]} 起，库默尔甚至能够利用狄利克雷的“抽屉方法”证明“理想因子”的“类”为数有限。\(^{5}\)

就类数的确定以及费马定理的证明在各种情形下的应用而言，我们将不去追述库默尔此后关于分圆域的工作的历史。这里只提一下他在 1859 年如何推广其方法，以便（至少部分地）在“库默尔域”\(\mathbf{Q}(\zeta, \mu)\) 中得到“素理想数”，其中 \(\mu\) 是一个不可约多项式 \(P(X) = X^p - \alpha\) 的根，而 \(\alpha \in \mathbf{Z}[\zeta]\) {[}188 e{]}。有趣的是，库默尔恰恰是把这个问题看作：把 \(\mathbf{Q}(\zeta, \mu)\) 视为取 \(\mathbf{Q}(\zeta)\) 为“基域”时的循环扩张：\(^6\)他从一个“素理想数”\(q\)（属于 \(\mathbf{Z}[\zeta]\)）出发，假设它既不整除 \(p\) 也不整除 \(\alpha\)，而这次他考察（用现代的术语）多项式 \(\overline{P}(X) = X^p - \overline{\alpha}\) 在剩余域 \(k\)（即 \(\mathbf{Q}(\zeta)\) 的与 \(\mathfrak{q}\) 相应的那个赋值的剩余域）中的情形（\(\overline{\alpha}\) 是 \(\alpha\) 在 \(k\) 中的典范像）。由于 \(\mathbf{Q}(\zeta)\) 是 \(p\) 次单位根的域，\(\overline{P}\) 或者在 \(k\) 上不可约，或者是一些一次因子的乘积。在第一种情形，库默尔说 \(\mathfrak{q}\) 在 \(\mathbf{Z}[\zeta, \mu]\) 中保持为素；在第二种情形，他引入元素 \(w_i(1 \leq i \leq p)\)（属于 \(\mathbf{Z}[\zeta]\)），它们在 \(k\) 中的像是 \(\overline{P}\) 的根，并且与每个指标 \(i\) 相对应地，他联系一个素理想因子 \(\mathfrak{r}_i\)（\(\mathfrak{q}\) 的因子）；接着令 \(W_i(X) = \prod_{j \neq i}(X - w_j)\)，他说，对于一个多项式 \(f\)（系数在 \(\mathbf{Z}[\zeta]\) 中），\(f(\mu)\) 含有理想因子 \(\mathfrak{r}_i\)；含 \(m\) 次，如果有

\[
f (w _ {i}) W _ {i} ^ {m} (w _ {i}) \equiv 0 \pmod {\mathfrak {q} ^ {m}}
\]

但

\[
f (w _ {i}) W _ {i} ^ {m + 1} (w _ {i}) \not \equiv 0 \pmod {q ^ {m + 1}}.
\]

总而言之，他就这样得到了 \(\mathbf{Q}(\zeta,\mu)\) 的在 Q 上非分歧的赋值，这对他心目中的应用而言已经足够。

库默尔在研究中很幸运：在研究因费马定理方面的工作最初把他引向的那些特殊的域时，他遇到了许多使研究变得容易得多的偶然情况。而库默尔的结果向一般情形的推广则呈现出巨大的困难，并要付出多年的努力。

在库默尔的发现之后的 40 年里，代数数理论的历史由担任主角的克罗内克与戴德金上演，它并非不令人想起（所幸没有同样尖刻的色彩）180 年前牛顿与莱布尼茨围绕无穷小演算的发明而展开的竞争。克罗内克是库默尔在柏林的学生，不久又成为其同事（如前所见，他的学位论文在库默尔理论的一个本质要点上帮了库默尔的忙），他对“理想数”抱有极其密切的兴趣，为的是把它们应用于他自己的研究；当看到他早在 1853 年（{[}186 a{]}，第 IV 卷，第 10 页）就陈述了关于 Q 的阿贝尔扩张的结构的一般定理，而且——或许更为惊人的是——在随后几年里创立了复乘理论并发现了类域论的最初萌芽（{[}186 a{]}，第 IV 卷，第 177-183 页与 207-217 页）时，我们对他的惊人洞察力深表钦佩。克罗内克 1857 年致狄利克雷的一封信（{[}186 a{]}，第 V 卷，第 418-421 页）表明，他在当时就已经掌握了库默尔理论的一个推广，这一点此外还由库默尔本人在他自己的一篇著作中予以证实（{[}188 e{]}，第 57 页），而且在 1860 至 1880 年间，克罗内克将在他的论文中经常援引这个理论。\(^{7}\)

但是，尽管此时德意志数论学派的数学家无人不晓克罗内克的这项工作，后者似乎只向一个有限的朋友与学生的圈子传达过其方法原理；而当他最终决定把它们发表出来——在 1881 年关于判别式的论文（{[}186 a{]}，第 II 卷，第 193-236 页）中，尤其是在 1882 年他宏大的“Festschrift”（纪念文集）（{[}186 a{]}，第 II 卷，第 237-287 页）中——时，戴德金不禁表达了他的惊讶（{[}79{]}，第 III 卷，第 427 页），因为他从听到的风声设想出的是相当不同的做法（{[}79{]}，第 III 卷，第 287 页）。此外，克罗内克远不具备戴德金那样程度的卓越阐述与明晰的天赋，因此毫不奇怪，主要是由后者（早在 1871 年就已发表）的方法构成了代数数论的框架；克罗内克的“未定元的添加”方法固然有趣，但就数论而言，在我们看来它不过是戴德金方法的一个变体；主要是在另一个指向代数几何的方向上，克罗内克的思想才获得了其对交换代数历史的全部重要性，这一点我们稍后将会看到。

出于要到很久以后才能清楚显现的理由，对任何一般理论的尝试而言，第一道序曲当然是澄清代数整数的概念。这一点大约在 1845 至 1850 年间获得，尽管很难为其出现标出一个确切的日期；看起来很可能是这样一个想法——一个对加法与乘法封闭的系统（或者更确切地说，我们现在所说的有限秩 Z-代数）——或多或少自觉地导致了代数整数的一般定义：当人们强使 Z{[}θ{]} 型的 Z-代数具有有限秩时（这是类比于由一个单位根生成的环 Z{[}ζ{]}，后者正处于当时算术学家们关切的中心），这个定义是不可避免的。因此，当狄利克雷（{[}92{]}，第 I 卷，第 640 页）、埃尔米特（{[}159{]}，第 I 卷，第 115 页与第 146 页）和艾森斯坦（{[}102 c{]}，第 236 页）彼此独立地引入代数整数的概念时，他们似乎既不认为这构成一个新观念，也不认为值得对它作一项详细研究；只有艾森斯坦实质上证明了（同前引文）两个代数整数的和与积仍是代数整数，况且他并未声称这一结果是独创的。

一个隐藏得深得多的问题，是确定可以在其中希望推广库默尔理论的那些环。库默尔本人在他的第一篇短文 {[}188 b{]} 中毫不犹豫地断言，通过考虑环 \(Z[\sqrt{D}]\)（D 为整数），他能用他的方法重新得到高斯的二元二次型理论；他从未发展过这个想法，但可以肯定的是，无论他还是戴德金之前的任何人，都未能看出：在形如 \(Z[\sqrt{D}]\) 的环中，当 \(D \equiv 1 \pmod{4}\) 时，素“理想”因子的唯一分解是不可能的（尽管三次单位根的例子已经表明，自高斯以来就被考虑的环 \(Z[\rho]\) 不同于 \(Z[\sqrt{-3}]\)）。⁸在戴德金和克罗内克之前，人们研究的环仍然只有 \(Z[\theta]\) 型的，或者有时是某些特定的 \(Z[\theta, \theta']\) 型的环。⁹就克罗内克而言，考虑一个代数扩张中全体整数的环这一想法，有可能最初是由代数函数域的研究提示给他的，在这个领域中，该环自然而然地表现为“在有限距离处取有限值”的函数全体所成之集：无论如何，他在 1881 年关于判别式的论文（早在 1862 年即已写成并向柏林科学院宣布）中，着重指出了这些域中“整数”的这一刻画（{[}186 a{]}，第 II 卷，第 193-236 页）。戴德金没有就他自己在这一点上的想法来源给出任何提示，但从 1871 年他关于数域的第一批出版物起，这样一个域的全体整数所组成的环便在他的理论中起着根本的作用；也正是戴德金，通过引入导子的概念，澄清了这样一个环与具有同一分式域的诸子环之间的关系（{[}79{]}，第 I 卷，第 105-157 页）。

但这并非唯一的困难。为了推广库默尔的想法，首先必须避免经过分解域，在非阿贝尔域的情形，它自然不可能有任何类似物。这个迂回乍看起来无论如何都非常令人惊讶而且矫揉造作，因为如果从不可约多项式 \(\Phi_p(X)\)（属于 \(\mathbf{Z}[X]\)）出发，人们会自问，为什么库默尔不把其想法的逻辑后果追究到底，是什么东西阻止他使用当时众所周知的“伽罗瓦虚数”理论。这个障碍在塞林——戴德金的一名学生——于 1865 年就已作出的一次不成功的推广尝试的映照下显得更清楚了：给定一个不可约多项式 \(P \in \mathbf{Z}[X]\)，塞林把相应的多项式 \(\overline{P}(X)\) 在 \(\mathbf{F}_q[X]\) 中分解成不可约因子；于是这个多项式的根属于一个有限扩张 \(\mathbf{F}_r\)（它是 \(\mathbf{F}_q\) 的扩张）；但塞林为了按库默尔同样的方式定义 \(q\) 的一个“素理想因子”在 \(P(X)\) 的根域中一个整数里的指数，毫不犹豫地在域 \(\mathbf{F}_r\) 中谈论模 \(q\) 的幂的同余（{[}282{]}，第 26 页）；而稍后，当他试图处理分歧问题时，他又向 \(\mathbf{F}_r\) “添加”形如 \(x^h = q\) 的方程的“虚根”（{[}282{]}，第 34 页）。显然，这些大胆的步骤（若把有限域 \(\mathbf{F}_q\) 换成 \(q\) 进域，它们本可得到辩护）在当时不可能导致别的结果，只能导致荒谬。所幸戴德金刚刚在 1857 年（{[}79{]}，第 I 卷，第 40-66 页）以“高阶同余理论”之名、用另一种形式重新研究了有限域的理论：\(^{10}\)他把这些域的元素解释为 Z{[}X{]} 中多项式关于一个“双重模”的“剩余”，该双重模由线性组合构成——系数取自 Z{[}X{]}——这些组合是素数 p 与一个首一不可约多项式 \(P \in Z[X]\) 的组合（这无疑正是他、也同样是克罗内克走向一般模概念的开端，稍后他们将各自独立地达到这一概念）。根据他本人的证词（{[}79{]}，第 I 卷，第 218 页）可以看出，戴德金曾着手处理 p 在域 \(Q(\xi)\) 中的“理想因子”问题（其中 \(P \in Z[X]\) 是 \(\xi\) 的极小多项式），方式如下（至少在“非分歧”的情形，也就是说，当 \(F_p[X]\) 中与 P 相应的多项式 \(\overline{P}\) 没有重根时）：我们在 Z{[}X{]} 中写出

\[
P = P _ {1} P _ {2} \dots P _ {h} + p G
\]

其中 \(\overline{P}_i\) 在 \(\mathbf{F}_p[X]\) 中不可约且互异；可以假设 \(G\)（在 \(\mathbf{Z}[X]\) 中）不被任何 \(P_i\) 整除，并且对所有 \(i\)，令 \(W_i = \prod_{j \neq i} P_j\)；于是，若 \(f \in \mathbf{Z}[X]\)，当下列条件成立时，我们就说 \(f(\xi)\) 含“理想因子”\(\mathfrak{p}_i\)（即 \(p\) 的对应于 \(P_i\) 的理想因子）\(k\) 次，如果有

\[
f W _ {i} ^ {k} \equiv 0 \pmod {p ^ {k}, P}
\]

且

\[
f W _ {i} ^ {k + 1} \not \equiv 0 \pmod {p ^ {k + 1}, P}.
\]

这里与库默尔对“库默尔域”所循方法的关系是显而易见的，而且人们同样可以用这个方式轻易地重新接上库默尔对分圆域的初始定义（例如可参见佐洛塔廖夫的著作 {[}343{]}，他起初独立于戴德金，稍后一点发展了这些想法）。

然而，无论是戴德金，还是似乎也作过类似尝试的克罗内克，都没有沿着这条道路继续走下去，两人都被分歧所呈现的困难挡住了（{[}79{]}，第 I 卷，第 218 页与 {[}186 a{]}，第 II 卷，第 325 页）。\(^{11}\)如果认为数域 K 的整数环 A 具有一个由同一个整数 \(\theta\) 的幂组成的（Z 上的）基，那么把上述方法推广到 Z{[} \(\theta\) {]} 中的分歧素数并无困难（正如佐洛塔廖夫所指出的（同前引文））。但存在这样的域 K：环 A 中不存在这种类型的基；而戴德金最终甚至发现，存在这样的情形：某些素数 p（域 K 的判别式的“异常因子”）具有这样的性质：不论 \(\theta \in A\) 如何选取，把上述方法应用于 \(\theta\) 在 \(\mathbf{Q}\) 上的极小多项式，都会导致把多重理想因子归于 \(p\)，而实际上 \(p\) 在 \(A\) 中并不分歧；\(^{12}\)他承认曾在这个未曾预见的困难面前停滞良久，之后才通过完整地创立模与理想的理论而成功地克服它，这一理论以精湛的方式（而且与其同时代人的散漫文风相对照，已经完全现代化）阐述于他无疑是其代表作的著作——狄利克雷数论著作的著名“第十一附录”（{[}79{]}，第 III 卷，第 1-222 页）——之中。这部著作将先后有三个版本，但从第一个版本（1871 年作为狄利克雷著作第二版的“第十附录”发表）起，方法的要点即已就位，而且几乎是单刀直入地，代数数理论从先前的草图与摸索跨入了一门完全成熟、已拥有其本质工具的学科：从一开始，数域的全体整数所组成的环就被置于理论的中心；戴德金证明了此环在 \(\mathbf{Z}\) 上基的存在性，并由此导出域的判别式的定义，即由整数环的一个基的元素及其共轭元组成的行列式的平方；不过在第十一附录中，他只对二次域给出了分歧素数（作为判别式的素因子）的一个刻画（{[}79{]}第 III 卷，第 202 页），尽管他自 1871 年起就已掌握一般定理。\(^{13}\)这部著作的中心结果是理想分解为素因子的存在性与唯一性定理，为此戴德金首先发展了一门初等的“模”理论；事实上，在第十一附录中，他把这个名称保留给数域的 Z-子模，但他所建立的概念以及他所证明的结果，其表述方式已经可以直接应用于最一般的模；\(^{14}\)其中我们必须特别指出，自 1871 年起引入的“诺特商”概念，它在唯一分解的第一个证明中起着重要的作用（“升链条件”同样如此）。在随后的两个版本中，戴德金又给出了这个定理的另外两个证明，他有理由把该定理视为这一理论的拱顶石。这里必须指出，正是在第三个证明中出现了分式理想（库默尔自 1859 年起已就分圆域引入）以及它们构成一个群这一事实；我们稍后将回到第二个证明（第 107 页）。

所有这些结果（除去所用的语言）无疑在 1860 年前后已为克罗内克所知，它们是他那些更一般的概念（我们稍后将谈到）的特殊情形（而戴德金承认，直到 1869-1870 年他才克服这一理论的最后困难（{[}79{]}，第 I 卷，第 351 页））；\(^{15}\)就数域而言，我们必须特别强调，从那时起，克罗内克就知道，当从一个本身为数域（异于 Q）的“基域”k 出发时，整个理论不作任何本质的改动即可适用，复乘理论曾自然地引出这一观点；他还对某些域 k 认识到了在 k 上非分歧的代数扩张 \(K \neq k\) 的存在性（{[}186 a{]}，第 II 卷，第 269 页），而这对于 k = Q 是不可能发生的（这可由埃尔米特与闵可夫斯基关于判别式的结果得出）。戴德金从未发展过最后这个观点（尽管他在 1882 年关于差积的论文中指出了其可能性），而“相对域”理论的第一个系统阐述应归功于希尔伯特（{[}163 a{]}，第 I 卷，第 63-363 页）。

最后，在 1882 年（{[}79{]}，第 I 卷，第 359-396 页），戴德金通过引入差积完成了这一理论，差积给了他判别式的一个新定义，并使他得以确定后者分解式中素理想因子的指数。也是在大约这个时候，他开始关注伽罗瓦扩张所引出的特殊问题，引入了分解群与惰性群的概念（在一篇直到 1894 年才发表的论文（{[}79{]}，第 II 卷，第 43-48 页）中），甚至（在生前未发表的文章（{[}79{]}，第 II 卷，第 410-411 页）中）给出了分歧群的一个概要，希尔伯特（独立于戴德金）稍后将把它发展开来（{[}163 a{]}，第 I 卷，第 13-23 页与 63-363 页）。

于是，大约在 1895 年，代数数论已结束其发展的第一阶段；在这一形成时期锻造出来的工具，将使它几乎可以立即着手下一阶段，即类域论的一般理论（或者说一回事，数域的阿贝尔扩张理论），这一理论一直延续至今，而我们在这里无须对其加以描述。从交换代数的观点看，可以说与此同时，戴德金环的理论除了其公理化刻画之外已实际上完成，这些环上有限型模的结构亦然（在数域的情形，它直到 1912 年才由施泰尼茨在实质上加以澄清 {[}294 b{]}）。\(^{16}\)

交换代数后来的进展将主要源自一些十分不同的问题，它们出自代数几何（此外，代数几何甚至在当代的“抽象”发展之前，就已经直接影响了数论）。

我们不打算在这里详述代数几何的历史，直到黎曼去世为止，它几乎没有触及我们的主题。只需记住，它首先以研究复射影平面中的代数曲线为目标，且大多用射影几何的方法（用或不用坐标）来处理。与此并行，在阿贝尔、雅可比、魏尔斯特拉斯和黎曼那里，单复变“代数函数”及其积分的理论发展起来；这一理论与平面代数曲线几何之间的联系在他们心目中是显然的，而且人们那时已经懂得如何“把分析应用于几何”；但用来研究代数函数的方法主要是“超越”性的，甚至在黎曼之前即已如此；\(^{17}\)这一特征在后者（黎曼）的工作中更为显著，那里引入了“黎曼曲面”以及定义在这种曲面上的任意解析函数。黎曼去世后几乎立即，罗赫特别是克莱布什就认识到，可以从黎曼用超越方法得到的深刻结果中，为曲线的射影几何提取出众多而惊人的应用，这自然会鼓励当时的几何学家为这些结果给出纯“几何”的证明；这一纲领由克莱布什和戈丹不完全地遵循，几年后由布里尔和 M. 诺特借助对给定曲线上变动点组以及经过这些点组的辅助曲线（“伴随曲线”）的研究而完成 {[}37{]}。但即使对同时代人来说，黎曼的超越方法（特别是他对拓扑概念和“狄利克雷原理”的使用）似乎也建立在不确定的基础之上；即使布里尔和诺特比大多数同时代的“综合”几何学家更谨慎（见后文第 106 页），他们的几何-解析论证也不是无可指摘的。正是为了给平面代数曲线理论一个坚实的基础，戴德金和韦伯才于 1882 年发表了他们关于这一主题的伟大著作 {[}80{]}：“在此发表的研究”，他们说，“其目标是：以一种同时简单、严格而且完全一般的方式，奠定单变量代数函数理论——黎曼的主要创造之一——的基础。在此后关于这一主题的研究中，常常对所考虑函数的奇异性加上限制性假设，而所设想的例外情形，或者作为极限情形顺带提及，或者完全被忽略。同样，某些关于连续性或解析性的基本定理是被假定的，它们的“自明性”依赖于各种性质的几何直观”（{[}80{]}，第 181 页）。\(^{18}\)

他们这项工作的本质想法，是把单变量代数函数的理论照搬到戴德金刚刚发展起来的代数数理论上；为此，他们必须首先采取一种“仿射”的观点（与其同时代人相反，后者无一例外地把代数曲线看作嵌入复射影空间中的对象）：因此，他们从有理分式域 \(\mathbf{C}(X)\) 的一个有限代数扩张 K 出发，并从 K 中“整代数函数”所组成的环 A 出发，即由该域中在多项式环 \(\mathbf{C}[X]\) 上为整的那些元素组成；他们的基本结果——不使用任何拓扑考虑而得到\(^{19}\)——是 A 为一个戴德金环，从而“第十一附录”的全部结果经必要的修改后都对它适用（甚至，如戴德金和韦伯虽然尚未看清原因却已指出的那样（{[}79{]}，第 I 卷，第 268 页），以更简单的方式适用）。此后，他们证明其定理实际上是双有理不变的（换句话说，只依赖于域 K），特别是不依赖于开头所作的“无穷远直线”的选取。对我们来说无疑更有意思的是：为了定义与 K 相应的“黎曼曲面”的点（特别是“无穷远点”，它们不可能对应 A 的理想），他们被引向引入域 K 的“位”的概念：他们面临的处境，正是盖尔范德 1940 年建立赋范代数理论时将再次遇到的处境，即有一个由元素组成的集合 K，这些元素并未预先作为函数给出，然而人们仍希望能够把它们当作函数来考虑；而为了得到这些假设中的函数的定义集合，他们第一次有了这样的想法（盖尔范德将再次采用、并且在现代数学中因被处处使用而变得平淡无奇的想法）：把集合 E 的一点 x、以及 E 到某个集合 G 的映射所成之集 F，与映射 \(f \mapsto f(x)\)（\(\mathcal{F}\) 到 \(G\) 的映射）联系起来，换句话说，在表达式 \(f(x)\) 中把 \(f\) 看作变量而把 \(x\) 看作固定的，这与经典传统相反。最后，从位的概念出发，他们毫不困难地定义了“正除子”（他们的术语是“Polygone”），它细化了 \(A\) 的理想的概念，并与布里尔和诺特的“点系”相对应；但是，虽然他们把主除子和微分的除子写成正除子的“商”，他们却没有给出除子的一般定义，直到 1902 年，亨塞尔与兰茨贝格才类比于分式理想而引入这一概念，它将一直使纯“几何”方法的拥护者们感到为难（他们不得不违心地以“虚系统”之名来定义它，却又无法给它一个“具体”的解释）。

同一年 1882 年还见证了人们等待了二十多年的克罗内克的伟大著作的问世（{[}186 a{]}，第 II 卷，第 237-387 页）。它比戴德金-韦伯的工作野心大得多，不幸的是也更加含糊和晦涩。它的中心主题是（用现代语言说）研究多项式环 \(\mathbf{C}[X_1,\ldots ,X_n]\) 或 \(\mathbf{Z}[X_1,\ldots ,X_n]\) 之一上的一个有限整代数的理想；克罗内克先验地把自己限于这些理想中为有限型的那些（它们全属此列这一事实，只是若干年后才由希尔伯特在其关于不变量的工作中、就 \(\mathbf{C}[X_1,\ldots ,X_n]\) 的理想加以证明（{[}163 a{]}，第 II 卷，第 199-257 页））。就 \(\mathbf{C}[X_1,\ldots ,X_n]\) 或 \(\mathbf{Z}[X_1,\ldots ,X_n]\) 而言，人们自然会把由一个理想的全部元素的公共零点所组成的“代数簇”与这两个环中每一个的理想相联系；而十九世纪期间所作的二、三维几何研究，将直观地引出这样的想法：每个簇都是数目有限的“不可约”簇之并，这些簇的“维数”未必相同。看来，这个事实的证明正是克罗内克为自己设定的目标，尽管他没有在任何地方明确陈述这一点，而且在他的著作中任何地方都找不到“不可约簇”或“维数”的定义。实际上，他只限于概括地指出，一个一般的消元方法如何从所考虑理想的生成元组出发，给出数目有限的代数簇，对其中每个簇，在适当的坐标系中，若干个坐标是任意的，而其余坐标是“代数函数”。\(^{21}\)但是，如果克罗内克瞄准的确实是分解为不可约簇，那就必须承认，他只在主理想这一初等情形取得了成功；在这一情形，他实质上通过推广高斯关于 Z{[}X{]} 的一个经典引理（{[}124 a{]}，第 I 卷，第 34 页），证明了环 C{[}X₁, \ldots, Xₙ{]} 与 Z{[}X₁, \ldots, Xₙ{]} 是唯一分解环；而在一般情形，甚至可以问，克罗内克当时是否已掌握素理想的概念（他所谓的“Primmodulsystem”是指不能分解为另外两个理想之乘积的理想（{[}186 a{]}，第 II 卷，第 336 页）；这一点尤其令人惊讶，因为戴德金自 1871 年以来给出的定义是完全一般的）。不过，稍后克罗内克改正了这个疏失。

不过必须指出，克罗内克的消元方法若运用得当，恰好导向代数簇分解为其不可约分支：这正是 E. 拉斯科尔在其 1905 年关于多项式理想的伟大著作开头清楚建立的事实 {[}194{]}；他正确定义了不可约簇的概念（在 \(\mathbf{C}^n\) 中）：代数簇 \(V\) 使得两个多项式的乘积不能在整个 \(V\) 上为零，除非其中之一在它上面为零；他还给出了一个不依赖于所选坐标轴的维数定义。在他插入这项工作的有趣的历史评述中，拉斯科尔还指出，他不仅接续了克罗内克和戴德金的纯代数倾向，而且接续了克莱布什和 M. 诺特学派的几何方法所提出的问题，特别是后者在 1873 年证明的著名定理 {[}239{]}。用今天的话说，这个定理实质上就是确定理想 \(\alpha\)——即 \(\mathbf{C}[X_1,\ldots ,X_n]\) 中那些在一个给定集合 \(M\)（位于 \(\mathbf{C}^n\) 中）的点上为零的多项式所组成的理想；大多数时候，\(M\) 是有限多个多项式 \(f_{i}\) 的公共零点的“代数簇”，而且长期以来人们似乎（当然没有论证）假设，至少对 \(n = 2\) 或 \(n = 3\)，理想 \(\alpha\) 就由 \(f_{i}\) 生成；\(^{22}\)M. 诺特已证明，即使对 \(n = 2\) 以及两个多项式 \(f_{1},f_{2}\)，这在一般情况下也是不正确的，而且他给出了使 \(\alpha\) 由 \(f_{1}\) 与 \(f_{2}\) 生成的充分条件。十年后，内托证明，不对 \(f_{1}\) 与 \(f_{2}\) 作任何假设，\(\alpha\) 的某个幂总包含在由 \(f_{1}\) 与 \(f_{2}\) 生成的理想中 {[}231{]}，希尔伯特 1893 年在其著名的“零点定理”中推广了这个定理（{[}163 a{]}，第 II 卷，第 287-344 页）。无疑正是在这个结果的启发下，拉斯科尔在其著作中引入了准素理想的一般概念\(^{23}\)，这是对环 \(C[X_{1},\ldots,X_{n}]\) 与 \(Z[X_{1},\ldots,X_{n}]\) 而言的（此前他在这两个环中通过转录戴德金的定义给出了素理想的定义），并证明了\(^{24}\)这两个环中每个理想都存在准素分解。\(^{25}\)他似乎没有关心过这种分解的唯一性问题；是麦考莱稍后 {[}211{]} 引入了“浸入”与“非浸入”准素理想的区分，并证明后者以唯一方式确定，而前者不然。最后必须指出，拉斯科尔还依靠魏尔斯特拉斯的“预备定理”把他的结果推广到了在一点邻域内收敛的整幂级数环。他著作的这一部分无疑是在纯代数观点下考虑这个环的第一处，而拉斯科尔为此发展的方法将在 1938 年深刻影响克鲁尔，那年克鲁尔创立了局部环的一般理论（参见 {[}187 h{]}，第 204 页及以下诸页）。

导向现代交换代数的观念运动，大约在 1910 年前后开始成形。如果说域的一般概念在二十世纪初已经获得，那么第一部定义了环的一般概念的工作，无疑要算弗伦克尔 1914 年的工作 {[}113 a{]}。那时，作为环的例子已经不仅有数论和代数

几何的整环，还有级数环（形式的或收敛的），最后还有基域上的（交换或非交换的）代数。然而，无论对环的理论还是对域的理论而言，起催化作用的似乎都是亨塞尔的 p 进数理论，弗伦克尔和施泰尼茨 {[}294 a{]} 都特别提到它是他们研究的出发点。

亨塞尔关于这一主题的第一篇出版物可追溯到 1897 年；在其中，他从戴德金和韦伯所揭示的类比出发，即一个代数函数域 \(K\) 的黎曼曲面的点与一个数域 \(k\) 的素理想之间的类比；他提出要把“皮瑟展开”（自十九世纪中叶以来即为经典）搬到数论中，这种展开在 \(K\) 的黎曼曲面的任一点的邻域内，使我们能把每个元素 \(x \in K\) 表示成所考虑点处“一致化元”的幂的收敛级数（一个只有有限多个负指数项的级数）。亨塞尔以同样的思路证明，如果 \(p\) 是数域 \(k\) 中位于素数 \(p\) 之上的一个素理想，就可以对每个 \(x \in k\) 联系一个“p 进级数”，形如 \(\sum_{i} \alpha_{i} p^{i}\)（或 \(\sum_{i} \alpha_{i} p^{i/e}\)，当 \(p\) 在 \(p\) 之上分歧时），其中 \(\alpha_{i}\) 取自理想 \(p\) 的商域的一组给定代表；但他的最大独创性在于有这样的想法：即使这些“展开”不对应任何 \(k\) 的元素，也去考虑它们，这是类比于黎曼曲面上超越函数的整级数展开 {[}157 a{]}。

在他此后的整个研究生涯中，亨塞尔都将致力于逐步打磨和完善他的新演算；如果说他的做法在我们看来可能迟疑而笨重，那么不应忘记，至少在开始时，他还不能使用今日数学的任何拓扑或代数工具，而这些工具本可使他的任务变得轻松。在他的早期出版物中，他此外几乎从不谈及拓扑概念，而且概括说来，对他而言 p 进数环（p 为数域 k 的整数环 A 的一个素理想），用现代的术语说，就是环 \(A/p^{n}\) 当 n 趋于无穷时的射影极限，不过带有纯代数的意义；而为了建立这个环及其分式域的性质，他必须在每一步都使用多少有些费力的临时性论证（例如为了证明 p 进数构成一个整环）。在 p 进域中引入拓扑概念的想法，在 1905 年之前几乎没有在亨塞尔那里出现 {[}157 d{]}；直到 1907 年，在写完那本按照他的思想重述代数数论的书 {[}157 f{]}之后，他才达到绝对 p 进赋值的定义及其基本性质 {[}157 e{]}，从这些出发，他将能够通过照搬柯西的理论，发展出一套完整的“p 进分析”，并懂得如何把它富有成效地应用于数论（特别是 p 进指数和对数的使用），其重要性自那时以来从未停止增长。

亨塞尔从一开始就清楚地看到，他的理论通过允许问题的“局部化”、并工作在这样一个域中——在那里不仅整除性质是平凡的，而且得益于他早在 1902 年就已得出的一条基本引理 {[}157 c{]}，对“简化”多项式 mod \(p\) 后没有重根的多项式的研究，被归结为对有限域上多项式的研究——而给经典命题带来了简化。他从 1897 年起 {[}157 b{]} 就给出了这些简化的惊人例子，特别是在涉及判别式的问题上（尤其是对他若干年前给出的“非凡因子”存在性判别准则的一个简短证明）。但在很长一段时间里，\(p\) 进数似乎在同时代数学家中引起了很大的不信任；这种态度无疑是对太“抽象”的观念常有的态度，不过其作者略微过分的热情（在数学中热衷于新理论的人之间如此常见）也不是全然无理地招致了它。事实上，亨塞尔不满足于把他的理论富有成效地应用于代数数，他还像所有同时代人一样深受 \(e\) 与 \(\pi\) 的超越性证明的触动，也许又被同时加于数和函数的“超越”这一称谓所迷惑，竟然到了这样的地步：认为他的 \(p\) 进数与实超越数之间存在某种联系，并且一度相信他由此能得到 \(e\) 甚至 \(e^e\) 的超越性的一个简单证明（{[}157 d{]}，第 556 页）。\(^{26}\)

1910 年之后不久，局面随着下一代人的崛起而改变，他们受弗雷歇和 F. 里斯的拓扑思想、施泰尼茨的代数思想的影响，并且从一开始就被“抽象”所征服；他们将懂得如何使亨塞尔的工作变得可以吸收并各归其位。从 1913 年起，屈尔沙克 {[}190{]} 以一般的方式定义了赋值的概念，认识到阿基米德赋值的重要性（p 进赋值正是它的一个例子），（通过照搬实数情形中的证明）证明了一个域关于一个赋值的完备化的存在性，尤其是以一般的方式证明了一个赋值扩张到给定域的代数扩张上的可能性。但他没有看到，赋值的非阿基米德特征已经隐藏在素域之中；这一点由奥斯特洛夫斯基建立，确定域 Q 上全部赋值的工作也应归功于他，还有把带有非阿基米德赋值的域刻画为 C 的子域的那条基本定理 {[}242{]}。在 1920 到 1935 年间的这几年里，理论通过对未必离散的赋值的更细致研究而趋于完备，其中包括对转到代数扩张或超越扩张时出现的各种不同情形的考察（奥斯特洛夫斯基、多林、F. K. 施密特）；另一方面，1931 年，克鲁尔引入并研究了赋值的一般概念 {[}187 f{]}，它在随后的岁月里将被扎里斯基及其代数几何学派大量使用。\(^{27}\)我们还必须在此提及——尽管这超出了我们的任务范围——关于完备赋值域和完备局部环结构的更深入的研究，它们出自同一时期（哈塞-施密特、维特、泰希米勒、I. 科恩）。

上面提到的那部弗伦克尔的工作（第 107 页）只处理了一种非常特殊的环（只有唯一素理想的阿廷环，而且该素理想还被假定是主的）。如果排除施泰尼茨关于域的工作 {[}294 a{]}，关于一般交换环研究的第一批重要工作就是 E. 诺特关于理想理论的两部伟大著作：1921 年的那部 {[}236 a{]}致力于准素分解，它以最一般的水平重新接过并在许多点上补全了拉斯科尔和麦考莱的结果；还有 1927 年的那部，它以公理化方式刻画了戴德金环 {[}236 b{]}。正如施泰尼茨就域所表明的那样，在这两部著作中可以看到，少数几个抽象观念——比如不可约理想的概念、链条件以及整闭环的思想（如我们所见，后两者已由戴德金提出）——如何能够单独引出一般的结果，而在这些结果事先已知的那些情形中，它们看起来是与纯计算的结果难分难解地联系在一起的。

随着 E. 诺特的这些著作，加上稍后阿廷-范德瓦尔登关于除子理想的工作（{[}317 a{]}，第 II 卷，第 105-109 页）以及克鲁尔把这些理想与本性赋值联系起来的工作 {[}187 f{]}，一个世纪之前开始的关于理想分解的漫长研究就此完成，\(^{28}\)与此同时，现代交换代数也宣告诞生。

关于交换代数此后不可胜数的研究，最容易按照几条主要的指导性倾向来加以归类：

\begin{enumerate}
\def\labelenumi{\Alph{enumi})}
\tightlist
\item
  局部环与拓扑环。尽管以胚胎的形式包含在数论和代数几何此前的全部工作之中，局部化的一般思想解脱出来得非常缓慢。分式环的一般概念直到 1926 年才由 E. 诺特的学生 H. 格雷尔定义，而且只对整环给出 {[}137{]}；它向更一般环的推广直到 1944 年才由 C. 谢瓦莱就诺特环给出，并在 1948 年由乌兹科夫就一般情形给出。直到大约 1940 年，克鲁尔及其学派几乎是仅有的在一般论证中使用一个整环的局部环 \(A_{\mathfrak{p}}\) 之考虑的人；
\end{enumerate}

这些环直到 1940 年起才随着谢瓦莱和扎里斯基的工作开始在代数几何中明确出现。\(^{29}\)

局部环本身的一般研究直到 1938 年才随着克鲁尔的伟大著作 {[}187 h{]}开始。这项工作最重要的结果涉及维数理论和正则环，我们不应当在此谈论它们；但也正是在那里，第一次出现了一个任意诺特局部环的完备化，以及与一个局部环相伴的分次环的一个尚不完善的形式；\(^{30}\)后者直到大约 1948 年才由 P. 塞缪尔 {[}270{]}以及勒雷和 H. 嘉当在代数拓扑方面的研究独立地加以定义。克鲁尔在刚才引用的著作中几乎不使用拓扑语言；但从 1928 年起 {[}187 e{]}，他已证明，在一个诺特环 A 中，单独一个理想 a 的诸幂之交，是由那些 \(x \in A\)——它们满足 \(x(1 - a) = 0\)（对某个 \(a \in a\)）——组成的集合；由此容易推出，对 A 的每个理想 m，A 上的 m 进拓扑在一个理想 a 上诱导出 a 的 m 进拓扑；在 1938 年的著作中，克鲁尔通过证明在一个诺特局部环中每个理想都是闭的而补全了这个结果。这些定理不久之后由谢瓦莱推广到半局部诺特环，随后由扎里斯基推广到冠以其名字的那些环 {[}340 b{]}；把“线性紧性”引入拓扑环，以及半局部完备环结构的确定，也应归功于谢瓦莱 {[}62 c{]}。

\begin{enumerate}
\def\labelenumi{\Alph{enumi})}
\setcounter{enumi}{1}
\tightlist
\item
  从局部到整体。自魏尔斯特拉斯以来，就存在这样的习惯：把一个单变量的解析函数（特别是代数函数）与它在黎曼曲面上有定义的所有点处的“展开”之集相联系。在他的数论一书的引言中（{[}157 f{]}，第 V 页），亨塞尔以同样的方式把一个代数数域 k 的每个元素与它在 k 的所有赋值\(^{31}\)处的完备化中与它对应的元素所成之集相联系。可以说，在现代交换代数中，正是这个观点取代了把一个理想分解为素理想之积的公式（在某种意义上推广了库默尔的初始观点）。亨塞尔的评述隐含地相当于把 k 嵌入它的所有完备化之积中；谢瓦莱 1936 年用他的“伊代尔”理论明确地做到了这一点 {[}62 b{]}，它完善了普吕弗和冯·诺伊曼先前的类似想法（后两人仅限于把 k 嵌入它的 p 进完备化之积中）。\(^{32}\)虽然这多少超出了我们的任务范围，但在此重要的是要指出：得益于伊代尔群上一个合适的拓扑，局部紧群的全部技术（包括哈尔测度）由此能够以非常有效的方式应用于数论。
\end{enumerate}

在更一般的框架中，克鲁尔的定理 {[}187 f{]}——把一个整闭环刻画为赋值环之交（这相当于把所考虑的环嵌入诸赋值环之积）——常常使这些环的研究变得容易，尽管这个方法只在克鲁尔环的本性赋值的情形才真正好用。此外，在克鲁尔 {[}187 i{]}那里还经常可以找到（相当初等的）例子，涉及那种“从局部到整体”的方法：把证明整环 A 的一个性质，化为对 A 在其所有素理想处的“局部化”\(A_{p}\) 逐一验证这一性质；\(^{33}\)更近些时候，塞尔注意到这个方法对任意的交换环 A 都有效，它也适用于 A-模及其同态，而且往往只需在 A 的极大理想处“局部化”就够了：这个观点与关于“谱”以及定义在这些谱上的层的思想紧密相连（见下文第 115 页）。

\begin{enumerate}
\def\labelenumi{\Alph{enumi})}
\setcounter{enumi}{2}
\tightlist
\item
  整数与整闭包。我们已经看到，最初为数域引入的代数整数的概念，已由克罗内克和戴德金推广到代数函数域，尽管在这一情形它可能显得相当人为（不对应于一个射影概念）。E. 诺特 1927 年的著作，随后是克鲁尔从 1931 年起的工作，将表明这些概念对于最一般的环所能提供的益处。\(^{34}\)我们特别要归功于克鲁尔的，是素理想在整代数中的提升定理 {[}187 g{]}，以及戴德金-希尔伯特的分解群与惰性群理论的推广 {[}187 f{]}。至于 E. 诺特，正规化引理的一般表述\(^{35}\)（希尔伯特的零点定理等即由之得出），以及第一个一般判别法（克罗内克和戴德金经典论证的转写，它使人们能够断言一个整环的整闭包在这个环上是有限的），都应归功于她。
\end{enumerate}

最后，必须在此指出，整闭环概念之所以在现代具有重要性，其原因之一来自扎里斯基关于代数簇的研究；他实际上发现，“正规”簇（也就是说其局部环为整闭的那些簇）以特别令人惬意的性质而著称，尤其是它们没有“余维数为 1 的奇点”这一事实；后来人们又注意到，类似的现象对“解析空间”也成立。同样，“正规化”（也就是说，对一个簇的各个局部环取其整闭包这一运算）本身也已成为现代代数几何武库中的一件强大工具。

\begin{enumerate}
\def\labelenumi{\Alph{enumi})}
\setcounter{enumi}{3}
\tightlist
\item
  模的研究与同调代数的影响。E. 诺特和 W. 克鲁尔在代数方面工作的标志性特征之一，是“线性化”的倾向，它延伸了戴德金和施泰尼茨铭刻在域理论上的类似方向；换句话说，理想首先是作为模来考虑的，因此人们被引向对它们应用线性代数的全部构造（商、积，以及更近些时候的张量积和同态模的构成），其结果一般地是已不再是理想的模。人们相当快地就看出，在许多问题中（无论涉及的是交换环还是非交换环），把自己局限于一个环 \(A\) 的理想的研究并无益处，相反，有必要就 \(A\)-模（最终还可加上某些有限性条件）更一般地陈述定理。
\end{enumerate}

同调代数的介入只是加强了先前的倾向，因为代数的这一分支本质上处理的是线性性质的问题。我们不打算在此重述它的历史；但值得指出的是，同调代数的几个基本概念（比如投射模的概念和函子 Tor 的概念）乃是对戴德金环上的模关于张量积的行为的细致研究的产物，这项研究是 H. 嘉当在 1948 年进行的。

反过来，可以预见，由同调代数作为 Ext 函子（投射模与内射模）和 Tor 函子（平坦模）的“泛零化子”而以自然方式引入的那些新模类，将给交换代数投去一道新的光。结果弄清楚的是，特别有用的是投射模，而更有用的是平坦模：后者的重要性首先来自塞尔 {[}283 b{]}所作的评注，即局部化和完备化自然地引入平坦模，从而以令人满意得多的方式“解释”了这两种运算已经为人所知的性质，并使它们的使用变得容易得多。此外还应当提及，同调代数的应用远不限于这些，它在代数几何中扮演着越来越深刻的角色。

\begin{enumerate}
\def\labelenumi{\Alph{enumi})}
\setcounter{enumi}{4}
\tightlist
\item
  谱的概念。按年代次序最后出现的交换代数新概念有着一段复杂的历史。希尔伯特的谱定理引入了希尔伯特空间中两两正交的投影算子所组成的有序集，它们构成一个“布尔代数”（或者更准确地说布尔格），\(^{36}\)并与 R 的（关于一个适当测度的）可测子集类的一个布尔格双射对应。大概正是他先前在希尔伯特空间算子方面的工作，使 M. H. 斯通在大约 1935 年去一般地研究布尔格，特别是去寻找它们用集合的子集（或某个等价关系之下的子集类）给出的“表示”。他观察到，当乘法由 \(xy = \inf(x, y)\) 定义、加法由 \(x + y = \sup(\inf(x, y'), \inf(x', y))\)定义时，一个布尔格就成了一个交换环（此外还是一个非常特殊类型的环）。在取有限集 X 的全部子集所组成的布尔格 \(\mathfrak{P}(X)\) 这一特殊情形中，立即可见 X 的元素与相应的“布尔”环的极大理想之间有一一的自然对应；而斯通恰好由此得到了他的一般定理：用包含一个元素的极大理想所成之集来表示布尔格 {[}301 a{]}。
\end{enumerate}

另一方面，人们早已知道，布尔格的一个经典例子是拓扑空间中既开又闭的子集所构成的集合。在第二项工作 {[}301 b{]}中，斯通证明，事实上每个布尔格也都同构于这种类型的一个布尔格。为此，自然必须在“布尔”环的极大理想所成之集上定义一个拓扑；这可以非常简单地做到：对每个理想 \(\mathfrak{a}\)，把包含 \(\mathfrak{a}\) 的极大理想所成之集取作闭集即可。

我们不打算在此谈论这些思想在泛函分析中的影响，在泛函分析中，它们在 I. 盖尔范德及其学派所发展的赋范代数理论的诞生中扮演了重要角色。但 1945 年，雅各布森观察到 {[}172 c{]}，斯通构想出的定义拓扑的程序，事实上可以应用到每个环 A（交换或非交换），只要取作理想之集的不是极大理想之集，而是双边“本原”理想之集（即使得 A/b 为本原环的那些双边理想 b）；对交换环，得到的当然是极大理想。就他而言，扎里斯基 1944 年 {[}340 a{]}用一个类似的方法在一个代数函数域的位所成之集上定义了一个拓扑。然而对大多数代数学家来说，这些拓扑仍然只是单纯的猎奇对象，因为它们一般不是豪斯多夫的，而且人们对使用如此笨拙的对象有一种相当可以理解的反感。这种不信任只是当 A. 韦伊在 1952 年表明每个代数簇都可以以自然的方式赋以刚才那种类型的拓扑、而且这个拓扑允许与微分簇或解析簇的情形完全类比地定义纤维空间的概念 {[}330 e{]}时，才得以消散；不久之后，塞尔有了把凝聚层理论推广到这样赋以拓扑的簇上的想法，多亏了这一理论，在“抽象”簇的情形，这个拓扑提供了与基域为 C 时通常的拓扑相同的服务，特别是在代数拓扑方法的应用方面 {[}283 a 与 b{]}。

从那时起，在全部交换代数中使用这种几何语言就是自然的了。人们很快观察到，极大理想的考虑对于得到方便的命题通常是不够的，\(^{37}\)恰当的概念是这个环的素理想所成之集，并以同样的方式赋以拓扑。随着谱概念的引入，我们现在手头有了一部词典，使得交换代数的每个定理都能用一种与韦伊-扎里斯基时期的代数几何十分接近的几何语言来表达；这还立刻导致了后者框架的相当可观的拓宽，以致从这一观点看，交换代数现在几乎只不过是其中最初等的部分 {[}138 a{]}。

